% ============================================================================
%  RESULTADOS DEL PLAN DE CALIDAD — ChibchaWeb
%  Documento compilable con pdflatex o lualatex
% ============================================================================
\documentclass[12pt,a4paper]{article}

% — Codificación y tipografía ------------------------------------------------
\usepackage[utf8]{inputenc}
\usepackage[T1]{fontenc}
\usepackage[spanish]{babel}
\usepackage{lmodern}

% — Diseño de página ---------------------------------------------------------
\usepackage[margin=2.5cm]{geometry}
\usepackage{parskip}

% — Tablas y figuras ----------------------------------------------------------
\usepackage{booktabs}
\usepackage{tabularx}
\usepackage{longtable}
\usepackage{multirow}
\usepackage{array}

% — Colores y resaltado -------------------------------------------------------
\usepackage[table,dvipsnames]{xcolor}
\definecolor{secbg}{HTML}{EBF5FB}

% — Hipervínculos y metadatos -------------------------------------------------
\usepackage{hyperref}
\hypersetup{
  colorlinks=true,
  linkcolor=MidnightBlue,
  urlcolor=MidnightBlue,
  pdfauthor={Equipo ChibchaWeb},
  pdftitle={Resultados del Plan de Calidad — ChibchaWeb},
}

% — Encabezados y pies de página ----------------------------------------------
\usepackage{fancyhdr}
\pagestyle{fancy}
\fancyhf{}
\fancyhead[L]{\small\textsc{ChibchaWeb}}
\fancyhead[R]{\small Resultados del Plan de Calidad}
\fancyfoot[C]{\thepage}
\renewcommand{\headrulewidth}{0.4pt}

% — Enumeración personalizada -------------------------------------------------
\usepackage{enumitem}

% ============================================================================
\begin{document}

% — Portada -------------------------------------------------------------------
\begin{titlepage}
\centering
\vspace*{2cm}
{\Huge\bfseries Resultados del Plan de Calidad\\[0.4em]
Proyecto ChibchaWeb\par}
\vspace{1.5cm}
{\Large\itshape Medición inicial de métricas de calidad\\tras el despliegue de la aplicación\par}
\vspace{2cm}
{\large
\begin{tabular}{rl}
\textbf{Proyecto:}   & ChibchaWeb \\
\textbf{Fecha:}      & 26 de mayo de 2026 \\
\end{tabular}\par}
\vspace{1.5cm}

{\large\textbf{Integrantes del equipo}\par}
\vspace{0.5cm}
\begin{tabular}{lll}
\toprule
\textbf{Nombre y apellidos} & \textbf{Código} & \textbf{Rol en el proyecto} \\
\midrule
Juan Carlos Duarte Sandoval & 20212020149 & Desarrollador \\
Sebastián Alejandro Bernate Guzmán & 20211020121 & Desarrollador \\
Laura Carina Álvarez Campos & 20212020006 & Product Owner \\
Johan Andrés Tamara Flautero & 20211020039 & Desarrollador \\
Laura Daniela Cubillos Escobar & 20211020045 & Product Manager \\
David Santiago Torres Mendieta & 20211020144 & Analista QA \\
\bottomrule
\end{tabular}

\vfill
{\small Universidad Distrital Francisco José de Caldas}
\end{titlepage}

% — Índice --------------------------------------------------------------------
\tableofcontents
\newpage

% ============================================================================
\section{Introducción}
% ============================================================================

Este documento tiene como propósito principal documentar los resultados iniciales del
proyecto ChibchaWeb, una aplicación integral de gestión de dominios. En un contexto
amplio, este proyecto se enmarca en la administración de perfiles de clientes, empleados,
distribuidores, procesamiento de pagos, registros de dominios y brindar soporte técnico.

ChibchaWeb es una plataforma esencial para la gestión y operatividad efectiva de los
servicios de dominio. Dado su alcance e impacto en las operaciones comerciales, la calidad
en su desarrollo, mantenimiento y operación es un aspecto crítico para el éxito global del
proyecto. El entorno de gestión de dominios es intrínsecamente complejo, abarcando desde la
interacción con usuarios finales hasta las complejidades técnicas de la gestión de registros. La
amplitud de esta aplicación abarca desde la autenticación y gestión de cuentas hasta la
implementación de transacciones financieras seguras y confiables.

A la fecha de publicación de este documento ya se ha realizado el despliegue de la aplicación,
por lo que se realiza una medición inicial de las métricas de calidad establecidos previamente
en el documento referente al Plan de Calidad en donde se presentan los siguientes objetivos de
calidad:

\begin{itemize}
  \item \textbf{Eficiencia del Sistema:} Garantizar un desempeño óptimo para las transacciones y
        consultas del sistema.
  \item \textbf{Disponibilidad:} Mantener el sistema disponible en un alto porcentaje del tiempo.
  \item \textbf{Seguridad de Datos:} Asegurar la integridad y confidencialidad de la información
        de usuarios y transacciones.
  \item \textbf{Satisfacción del Usuario:} Velar por una experiencia de usuario satisfactoria y
        fluida.
\end{itemize}

Y se espera obtener los siguientes resultados:

\begin{enumerate}
  \item \textbf{Tiempo de respuesta del sistema:} Se considerará exitoso si se mantiene por debajo
        de los 2 segundos para la mayoría de las operaciones.
  \item \textbf{Disponibilidad del sistema:} Se establece un objetivo de disponibilidad del 99,5\,\% en
        un período de un mes.
  \item \textbf{Seguridad de los datos:} Se medirá mediante auditorías de seguridad regulares,
        manteniendo un puntaje superior al 90\,\% en los estándares de seguridad.
  \item \textbf{Satisfacción del usuario:} Se espera un puntaje promedio de satisfacción superior a
        4,5 en una escala de 1 a 5.
\end{enumerate}

% ============================================================================
\section{Tiempo de respuesta del sistema}
% ============================================================================

Para medir el tiempo de respuesta del sistema se hizo uso de la herramienta de Google
PageSpeed (\url{https://pagespeed.web.dev}) y se obtuvieron los siguientes resultados.

Es posible consultar estos resultados en el siguiente enlace:
\url{https://pagespeed.web.dev/analysis/https-chibchaweb-pythonanywhere-com/1nwakfovb9?form_factor=desktop}

\subsection{Diagnóstico de problemas de rendimiento}

\begin{itemize}
  \item \textbf{Actuación (Performance):} Este indicador mide la velocidad de carga y rendimiento general
        del sitio web. Un puntaje alto en este ítem sugiere que el sitio web tiene una carga rápida y
        eficiente, lo que mejora la experiencia del usuario al navegar por las páginas.
  \item \textbf{Accesibilidad (Accessibility):} Evalúa la capacidad del sitio web para ser utilizado por
        personas con discapacidades o dificultades de acceso. Un puntaje alto indica que el sitio
        cumple con estándares de accesibilidad, facilitando su uso para una amplia gama de usuarios.
  \item \textbf{Buenas Prácticas (Best Practices):} Este ítem revisa si el sitio sigue las mejores prácticas de
        desarrollo web y seguridad en línea. Un puntaje perfecto sugiere que el sitio cumple con
        estándares y normativas de seguridad, optimización y desarrollo web.
  \item \textbf{SEO (Search Engine Optimization):} Evalúa la optimización del sitio para motores de
        búsqueda. Un puntaje alto sugiere que el sitio está bien estructurado y optimizado para ser
        encontrado fácilmente en los resultados de búsqueda.
\end{itemize}

Los puntajes altos en Actuación, Accesibilidad, Buenas Prácticas y SEO indican que el sitio
web está bien optimizado en términos de velocidad, accesibilidad, cumplimiento de
estándares de desarrollo y capacidad para ser indexado por los motores de búsqueda. Estos
puntajes reflejan un sólido rendimiento del sitio en diferentes aspectos, mejorando la
experiencia del usuario y la visibilidad en línea.

\subsection{Métricas}

\begin{itemize}
  \item \textbf{First Contentful Paint:} Es el tiempo que tarda el navegador en mostrar el primer contenido
        visual en la pantalla después de que se solicita la página. Un valor bajo indica que la página
        comienza a ser percibida rápidamente por los usuarios.
  \item \textbf{Renderizado del mayor elemento con contenido:} Mide el tiempo que tarda en cargar el
        elemento más grande o significativo de la página. Generalmente, se refiere a imágenes,
        videos o bloques de texto. Un valor bajo aquí indica que la mayor parte del contenido
        relevante se carga rápidamente.
  \item \textbf{Total Blocking Time:} Se refiere al tiempo total durante el cual la interactividad de la página
        está bloqueada para los usuarios debido a tareas de procesamiento de JavaScript. Un valor
        bajo es deseable ya que minimiza las interrupciones en la interacción del usuario.
  \item \textbf{Cambios de Diseño Acumulados:} Evalúa la estabilidad visual de la página durante la carga.
        Indica si los elementos de la página se mueven o cambian de posición mientras se carga, lo
        cual puede ser molesto para los usuarios. Un valor bajo significa menos cambios inesperados
        en el diseño.
  \item \textbf{Speed Index:} Es una métrica que cuantifica qué tan rápido se muestra el contenido visible
        durante la carga de la página. Un valor bajo indica que el contenido relevante se presenta
        rápidamente a los usuarios.
\end{itemize}

Los resultados reflejan una buena velocidad de carga inicial, una interacción sin bloqueos y
una estabilidad visual adecuada. Sin embargo, mejorar el tiempo de carga del elemento más
grande puede contribuir a una mejor experiencia de usuario al presentar más rápidamente el
contenido principal de la página.

\subsection{Diagnósticos}

Estos datos representan las oportunidades identificadas para mejorar el rendimiento y la
eficiencia de la página web. Cada diagnóstico está acompañado por una estimación del
ahorro potencial en tiempo de carga, medido en milisegundos, y una estimación de los bytes
consumidos por actividad medidos en KiB.

\begin{itemize}
  \item \textbf{Recursos que bloquean el renderizado de la página:} Se recomienda que se muestren los
        elementos de JavaScript y CSS críticos insertados y se pospongan todos los que no sean
        esenciales.
  \item \textbf{Reducir JavaScript no utilizado:} Optimizar el código JavaScript para eliminar el código que
        no se utiliza puede contribuir a una reducción de 62~KiB en la carga, eliminando recursos
        innecesarios.
  \item \textbf{Reducir CSS no utilizado:} Eliminar el código CSS que no se
        utiliza puede contribuir a ahorrar 58~KiB de carga, mejorando la eficiencia de la página.
\end{itemize}

Estas oportunidades de mejora, entre otras, ofrecen un potencial significativo para optimizar
el rendimiento de la página web, reduciendo considerablemente los tiempos de carga y
mejorando la experiencia general del usuario al acceder al sitio.

% ============================================================================
\section{Disponibilidad del sistema}
% ============================================================================

Para medir la disponibilidad del sistema se utilizó UptimeRobot, indicando una
disponibilidad del 100\,\%, también debido al poco tiempo de monitoreo y al host donde se
encuentra ubicada la página web, por lo que es necesario realizar más mediciones a futuro.

% ============================================================================
\section{Seguridad de los datos}
% ============================================================================

Para medir la seguridad de los datos se utilizó una herramienta gratuita llamada HostedScan.
Esta plataforma proporciona una evaluación detallada de diversas vulnerabilidades
potenciales presentes en el entorno del sitio web.

Los resultados obtenidos reflejan un entorno seguro, sin vulnerabilidades críticas ni de alto
riesgo. En total se detectaron 9 riesgos clasificados de la siguiente manera:

\begin{itemize}
  \item \textbf{Riesgos Críticos:} 0
  \item \textbf{Riesgos Altos:} 0
  \item \textbf{Riesgos Medios:} 2
  \item \textbf{Riesgos Bajos:} 7
  \item \textbf{Riesgos Aceptados o Cerrados:} 0
\end{itemize}

Entre las observaciones destacadas se encuentran configuraciones de seguridad HTTP no
establecidas, tales como la ausencia de la cabecera ``Strict-Transport-Security'' y la cabecera
``Content Security Policy (CSP)'', así como cookies sin el atributo ``SameSite''. Si bien estas
observaciones no representan una amenaza inmediata, deben ser tratadas con prioridad para
fortalecer la protección ante ataques del tipo \textit{man-in-the-middle}, \textit{cross-site scripting} (XSS)
o CSRF (\textit{Cross-Site Request Forgery}).

Este resultado evidencia que, si bien la plataforma no presenta vulnerabilidades severas, es
necesario continuar con la implementación de políticas de seguridad, actualización de
configuraciones de cabeceras HTTP, y fortalecimiento de las buenas prácticas de desarrollo
seguro.

% ============================================================================
\section{Satisfacción del usuario}
% ============================================================================

Para medir la satisfacción de los usuarios se realizó una encuesta en línea con 5 preguntas y se
les solicitó que indicaran en una escala del 1 al 5 qué tan satisfechos se encontraban con cada
uno de los ítems mencionados. Las preguntas realizadas y sus resultados son:

\begin{enumerate}
  \item \textbf{Experiencia de Navegación y Usabilidad:} Facilidad de uso para acceder y navegar
        por la plataforma. \textbf{Promedio: 5,0}
  \item \textbf{Rendimiento y Velocidad:} Rapidez en el acceso a las funcionalidades. \textbf{Promedio: 4,6}
  \item \textbf{Atención al Cliente:} Eficacia y calidad de la asistencia brindada por el equipo.
        \textbf{Promedio: 4,6}
  \item \textbf{Seguridad y Confidencialidad:} Percepción sobre la seguridad de los datos
        proporcionados. \textbf{Promedio: 4,8}
  \item \textbf{Cumplimiento de Expectativas:} En qué medida ChibchaWeb ha cumplido con las
        expectativas del usuario. \textbf{Promedio: 4,8}
\end{enumerate}

% ============================================================================
\section{Conclusiones}
% ============================================================================

\begin{itemize}
  \item Se obtuvo un tiempo de respuesta de 98/100, lo que indica una buena respuesta.
  \item En cuanto a la seguridad de los datos, no se tienen riesgos críticos o altos, por lo que se
        cumple con la métrica de calidad; sin embargo, hay que tener en cuenta los riesgos bajos y
        medios obtenidos para mejorar la seguridad de la plataforma.
  \item Se obtuvo una puntuación total promedio de 4,76 en cuanto a la satisfacción del usuario, lo
        que cumple con las métricas de calidad planteadas en un inicio.
  \item La disponibilidad es del 100\,\%, cumpliendo con lo esperado.
\end{itemize}

\end{document}
